\documentclass[10.5pt]{article}
\usepackage[letterpaper, margin=0.6in, top=0.5in, bottom=0.5in]{geometry}
\usepackage{amsmath,amssymb,amsthm}
\usepackage{fancyhdr}
\usepackage{xcolor}
\usepackage{tikz}
\usepackage{enumitem}
\usepackage{tcolorbox}
\usepackage{hyperref}

% Define colors matching instructor board annotations
\definecolor{chalkblue}{RGB}{0, 105, 185}
\definecolor{chalkred}{RGB}{205, 35, 35}
\definecolor{chalkpink}{RGB}{230, 90, 120}
\definecolor{chalkcyan}{RGB}{15, 160, 190}
\definecolor{chalkgreen}{RGB}{20, 130, 50}
\definecolor{boxbg}{RGB}{252, 252, 254}
\definecolor{boxframe}{RGB}{150, 160, 180}

\pagestyle{fancy}
\fancyhf{}
\renewcommand{\headrulewidth}{0pt}
\renewcommand{\footrulewidth}{0pt}

\setlength{\parindent}{0pt}
\setlength{\parskip}{2.5pt}

\begin{document}

% ==============================================================================
% PAGE 1
% ==============================================================================
\noindent \textbf{\large Math 105L} \hfill \textbf{\large Lecture 2-2}\\[1pt]
{\textbf{\Large Rational Functions: Intercepts, Asymptotes, \& Sign Charts}}\\[-5pt]
\rule{\textwidth}{1.2pt}

\vspace{0.03in}

\textbf{\large 1. Anatomy of a Rational Function}

A \textbf{rational function} is a quotient of two polynomials:
\[
  r(x) = \frac{p(x)}{q(x)} \quad \text{where } q(x) \ne 0
\]
\begin{tcolorbox}[colback=boxbg, colframe=boxframe, boxrule=0.8pt, arc=3pt, left=8pt, right=8pt, top=4pt, bottom=4pt]
\begin{itemize}[leftmargin=1.2em, topsep=1pt, itemsep=3pt]
  \item \textbf{$x$-Intercepts (Roots):} Occur where the numerator vanishes, $p(a) = 0$ (with $q(a) \ne 0$). The point is $(a, 0)$.
  \item \textbf{Vertical Asymptotes (VA):} Occur where the denominator vanishes, $q(b) = 0$ (with $p(b) \ne 0$). The line is $x = b$.
  \item \textbf{$y$-Intercept:} Evaluated at $x = 0$: $r(0) = \frac{p(0)}{q(0)}$. The point is $(0, r(0))$.
  \item \textbf{Horizontal Asymptotes (HA):} Determined by degrees of $p(x)$ and $q(x)$ as $x \to \pm\infty$. When degrees are equal, $y = \frac{\text{leading coefficient of } p}{\text{leading coefficient of } q}$.
\end{itemize}
\end{tcolorbox}

\vspace{0.03in}

\textbf{\large 2. Classroom Blackboard Example: $r(x) = \frac{(2x - 1)(x + 1)}{(x - 1)(x + 2)}$}

\begin{itemize}[leftmargin=1.2em, topsep=1pt, itemsep=3pt]
  \item \textbf{$x$-Intercepts:} Set numerator to zero:
  \[
    (2x - 1)(x + 1) = 0 \implies \mathbf{x = \frac{1}{2}}, \quad \mathbf{x = -1}
  \]
  \item \textbf{Vertical Asymptotes:} Set denominator to zero:
  \[
    (x - 1)(x + 2) = 0 \implies \mathbf{x = 1}, \quad \mathbf{x = -2}
  \]
  \item \textbf{$y$-Intercept:}
  \[
    r(0) = \frac{(2(0) - 1)(0 + 1)}{(0 - 1)(0 + 2)} = \frac{(-1)(1)}{(-1)(2)} = \mathbf{\frac{1}{2}} \implies \left(0, \frac{1}{2}\right)
  \]
  \item \textbf{Horizontal Asymptote:} $\lim_{x\to\pm\infty}\frac{2x^2 + x - 1}{x^2 + x - 2} = \frac{2}{1} = \mathbf{2} \implies \mathbf{y = 2}$.
\end{itemize}

\vspace{0.03in}

\textbf{\large 3. The Factor Sign Analysis Table}

The critical values $\{-2, -1, 1/2, 1\}$ partition the real line into five intervals:
\begin{center}
\begin{tabular}{|c|c|c|c|c||c|}
  \hline
  Interval & $(2x - 1)$ & $(x + 1)$ & $(x - 1)$ & $(x + 2)$ & \textbf{Total Sign of $r(x)$} \\
  \hline
  $x < -2$ & $-$ & $-$ & $-$ & $-$ & $\mathbf{+} \text{ (above axis)}$ \\
  $-2 < x < -1$ & $-$ & $-$ & $-$ & $+$ & $\mathbf{-} \text{ (below axis)}$ \\
  $-1 < x < 0.5$ & $-$ & $+$ & $-$ & $+$ & $\mathbf{+} \text{ (above axis)}$ \\
  $0.5 < x < 1$ & $+$ & $+$ & $-$ & $+$ & $\mathbf{-} \text{ (below axis)}$ \\
  $x > 1$ & $+$ & $+$ & $+$ & $+$ & $\mathbf{+} \text{ (above axis)}$ \\
  \hline
\end{tabular}
\end{center}

\vspace{0.02in}

\begin{tcolorbox}[colback=chalkblue!6, colframe=chalkblue, boxrule=1pt, arc=3pt, left=8pt, right=8pt, top=4pt, bottom=4pt]
\textbf{Sign Chart Rule:}\\[1pt]
The sign of $r(x)$ can only change at an $x$-intercept (crossing zero) or a vertical asymptote (blowing up across $\pm\infty$). Testing the sign of each linear factor within each partition completely determines where the rational curve lies above or below the $x$-axis.
\end{tcolorbox}

\end{document}

